\documentclass[../DoAn.tex]{subfiles}
\begin{document}
\section{Đặt vấn đề}
\label{section:1.1}
Đề tài tập trung vào nhu cầu tìm và đăng ký phòng học, phòng họp phục vụ giảng dạy và công việc chuyên môn. Trong kịch bản quản lý lịch bằng bảng tính mà nhóm lựa chọn, người phụ trách phải đối chiếu thủ công thời gian sử dụng, sức chứa và trang thiết bị trước khi xác nhận yêu cầu. Điều này làm chậm quá trình tìm phòng phù hợp và tăng nguy cơ trùng lịch.

Từ nhu cầu đó, nhóm phát triển RoomHub nhằm số hóa quy trình tìm phòng, đặt phòng, phê duyệt, thông báo và cấp quyền truy cập. Hệ thống hướng đến hai nhóm người dùng chính là giảng viên và cán bộ quản trị; hỗ trợ quản lý sức chứa, trang thiết bị, lịch sử dụng và hình thức khóa của từng phòng.

\section{Mục tiêu và phạm vi đề tài}
\label{section:1.2}
Phiên bản thử nghiệm ban đầu mô hình hóa một tòa nhà tám tầng và lấy tòa C7 của Đại học Bách khoa Hà Nội làm bối cảnh tham chiếu. Sơ đồ tầng và dữ liệu phòng do nhóm xây dựng, không mô tả mặt bằng thực tế hay việc triển khai trên toàn trường. Ứng dụng Android phục vụ giảng viên và cán bộ quản trị, bao gồm tìm và đặt phòng, phê duyệt yêu cầu, thông báo, quản lý phòng và lịch sử dụng, cùng cấp quyền tiếp cận theo hình thức khóa. Các thành phần máy chủ và lưu trữ được trình bày ở phần kiến trúc và giải pháp.

\section{Định hướng giải pháp}
\label{section:1.3}
Giải pháp gồm ứng dụng Android viết bằng React Native và TypeScript, cùng dịch vụ Java 17/Spring Boot. Dịch vụ cung cấp REST API và dùng Jakarta Persistence để lưu dữ liệu vào H2 ở chế độ tệp trong môi trường thử nghiệm. Cấu hình hiện tại chọn chế độ dùng máy chủ; chế độ lưu cục bộ được duy trì riêng và không tự hợp nhất dữ liệu phát sinh khi mất mạng. Đường truy cập công khai hiện chỉ phục vụ thử nghiệm, chưa bảo đảm tính sẵn sàng của một hệ thống vận hành thực tế. Các thao tác tạo và duyệt yêu cầu được kiểm tra trong giao dịch có khóa bi quan trên bản ghi cấu hình và phòng. Kiến trúc được trình bày tại Chương~\ref{chapter:Experiment}.

Đối với phòng gắn khóa số, việc duyệt yêu cầu cấp quyền tạo mã cho đúng người dùng và đúng phòng, không chuyển người dùng thành quản trị viên hệ thống. Mã tạm thời gồm bảy chữ số, có khoảng hiệu lực dựa trên lịch đặt phòng. Ứng dụng lưu mã trước và gửi bản tin MQTT trực tiếp qua OneIoT khi người dùng đã kết nối trong phiên. Việc khóa thật nhận và áp dụng mã là bước kiểm chứng riêng, không suy ra chỉ từ kết quả gửi bản tin.

\section{Bố cục báo cáo}
\label{section:1.4}
Chương~\ref{chapter:Related_works} phân tích bối cảnh, kịch bản nghiệp vụ, tác nhân, use case và quy tắc vận hành. Chương~\ref{chapter:Methodology} giải thích công nghệ và cơ sở lý thuyết được dùng. Chương~\ref{chapter:Experiment} trình bày thiết kế, kết quả triển khai, giao diện và kiểm thử. Chương~\ref{chapter:SolutionAndContribution} phân tích các giải pháp kỹ thuật nổi bật. Chương~\ref{chapter:conclusion} tổng kết kết quả và hướng phát triển.
\end{document}
